\documentclass[17pt]{beamer}

\usepackage[utf8]{inputenc}
\usepackage[turkish]{babel}
\usepackage[T1]{fontenc}

\usepackage{amsmath}
\usepackage{amssymb}
\usepackage{amsthm}
\usepackage{enumerate}
\usepackage{tikz}
\usepackage{transparent}
\usepackage{xcolor}

\usetheme{Warsaw}
\usecolortheme{default}
\definecolor{nigdeyesili_acik}{RGB}{3, 150, 166}
\definecolor{nigdeyesili_koyu}{RGB}{143, 209, 217}
\setbeamercolor{structure}{fg=nigdeyesili_acik}

\author[Dr. Mehmet CANEVİ]{Arş.~Gör.~Dr.~M.~Canevi\inst{1}}

\institute{
    \inst{1}%
    Bilgisayar Mühendisliği\\
    Mühendislik Fakültesi
}
    
\date[2025] {Ders Notları, Ocak 2025}
    
\logo{\transparent{0.4}\includegraphics[width = 20mm]{logo}}
    
\title[Ders 2] {Fonksiyonlar}
\begin{document}
%%%%%%%%%%%%%%%%%%%%%%%%%%%%%%%%%%%%%%%%%%%%%%%%%%%%%%%%%%%%%%%%%%%%%%%%%%%%%%%%
\frame{\titlepage}
\begin{frame}[fragile]{İçidekiler}
    \tableofcontents
\end{frame}
%%%%%%%%%%%%%%%%%%%%%%%%%%%%%%%%%%%%%%%%%%%%%%%%%%%%%%%%%%%%%%%%%%%%%%%%%%%%%%%%
\section{Fonksiyonları tanımlama}
\begin{frame}[fragile]{Fonksiyon nedir?}
    Tekrar kullanılmak üzere bölünmüş kod parçacıklarıdır
    \begin{lstlisting}
        veri_tipi fonksiyon_adi(veri_tipi arg1,veri_tipi arg1,...)
        {
            kod;
            return veri_tipinde_deger;
        }\end{lstlisting}
    veya
    \begin{lstlisting}
        void fonksiyon_adi(veri_tipi arg1,veri_tipi arg1,...)
        {
            kod;
        }\end{lstlisting}
    ile tanımlanır.
\end{frame}
%%%%%%%%%%%%%%%%%%%%%%%%%%%%%%%%%%%%%%%%%%%%%%%%%%%%%%%%%%%%%%%%%%%%%%%%%%%%%%%%
\begin{frame}[fragile]{Fonksiyon örneği}
    Matematikteki fonksiyonlara benzer $y=sin(x)$ gibi
    \lstinputlisting{lec2/sinx.c}
\end{frame}
%%%%%%%%%%%%%%%%%%%%%%%%%%%%%%%%%%%%%%%%%%%%%%%%%%%%%%%%%%%%%%%%%%%%%%%%%%%%%%%%
\begin{frame}[fragile]{main fonksiyonu}
    \lstinline{C} programları için 
    \lstinputlisting{lec2/main.c}
    fonksiyonu kullanılmalıdır. Döndürülen \verb|0| değeri işletim sistemine \verb|OK| döndürür. \verb|-1| ise \verb|ERROR| mesajıdır.
\end{frame}
%%%%%%%%%%%%%%%%%%%%%%%%%%%%%%%%%%%%%%%%%%%%%%%%%%%%%%%%%%%%%%%%%%%%%%%%%%%%%%%%
\begin{frame}[fragile]{Örnek}
    \lstinputlisting{lec2/main2.c}
\end{frame}
%%%%%%%%%%%%%%%%%%%%%%%%%%%%%%%%%%%%%%%%%%%%%%%%%%%%%%%%%%%%%%%%%%%%%%%%%%%%%%%%
\begin{frame}[fragile]{Toplama}
    \lstinputlisting{lec2/toplama.c}
\end{frame}
%%%%%%%%%%%%%%%%%%%%%%%%%%%%%%%%%%%%%%%%%%%%%%%%%%%%%%%%%%%%%%%%%%%%%%%%%%%%%%%%
\begin{frame}[fragile]{Değişkenin geçerliliği}
    \lstinputlisting{lec2/degisken1.c}
\end{frame}
%%%%%%%%%%%%%%%%%%%%%%%%%%%%%%%%%%%%%%%%%%%%%%%%%%%%%%%%%%%%%%%%%%%%%%%%%%%%%%%%
\begin{frame}[fragile]{Fonksiyon tanımı}
    \lstinputlisting{lec2/fonksiyon1.c}
\end{frame}
%%%%%%%%%%%%%%%%%%%%%%%%%%%%%%%%%%%%%%%%%%%%%%%%%%%%%%%%%%%%%%%%%%%%%%%%%%%%%%%%
\begin{frame}[fragile]{Matematik fonksiyonları}
    Bazı matematiksel fonksiyonlar:
    \begin{itemize}
    \item \lstinline{double sin(double);}
    \item \lstinline{double tan(double);}
    \item \lstinline{double fabs(double);}
    \item \lstinline{double log10(double);}
    \item \lstinline{double sqrt(double);}
    \end{itemize}
\end{frame}
%%%%%%%%%%%%%%%%%%%%%%%%%%%%%%%%%%%%%%%%%%%%%%%%%%%%%%%%%%%%%%%%%%%%%%%%%%%%%%%%
\begin{frame}[fragile]{Matematik fonksiyonları örneği}
    Bir kullanım örneği:
    \lstinputlisting{lec2/mat1.c}
\end{frame}
%%%%%%%%%%%%%%%%%%%%%%%%%%%%%%%%%%%%%%%%%%%%%%%%%%%%%%%%%%%%%%%%%%%%%%%%%%%%%%%%
\section{Sorular}
\begin{frame}[fragile]{Sorular}
    \begin{alertblock}{Soru-1}
        İki sayı ile 4 işlem yapan programı yazınız.
    \end{alertblock}
    % \begin{alertblock}{Soru-2}
    %     İki sayı ile 4 işlem yapan programı yazınız.
    % \end{alertblock}
\end{frame}
%%%%%%%%%%%%%%%%%%%%%%%%%%%%%%%%%%%%%%%%%%%%%%%%%%%%%%%%%%%%%%%%%%%%%%%%%%%%%%%%
\end{document}